    \def\stat{bistrov}





\def\tit{ОНТОЛОГИЯ ПОДДЕРЖКИ НЕПРЕРЫВНОСТИ ИТ-УСЛУГ


В~ИНФОРМАЦИОННО-ТЕЛЕКОММУНИКАЦИОННЫХ СИСТЕМАХ}





\def\titkol{Онтология поддержки непрерывности ИТ-услуг


в~ИТС} %информационно-телекоммуникационных системах}





\def\autkol{И.\,И. Быстров, В.\,Н.~Веселов, С.\,И.~Радоманов}





\def\aut{И.\,И. Быстров$^1$, В.\,Н.~Веселов$^2$, С.\,И.~Радоманов$^3$}





\titel{\tit}{\aut}{\autkol}{\titkol}





%{\renewcommand{\thefootnote}{\fnsymbol{footnote}}


%\footnotetext[1]


%{Работа выполнена при частичной финансовой поддержке  программы


%<<Интеллектуальные информационные технологии, системный анализ и


%автоматизация>> (проект 1.7).}}











\renewcommand{\thefootnote}{\arabic{footnote}}


\footnotetext[1]{Институт проблем информатики Российской академии наук, ibystrov@ipiran.ru}


\footnotetext[2]{Институт проблем информатики Российской академии наук, vveselov@ipiran.ru}


\footnotetext[3]{Институт проблем информатики Российской академии наук, radomanov@ipiran.ru}








\label{st\stat}





                  \thispagestyle{headings}





\vspace*{-6pt}





  \Abst{Рассматриваются возможности использования онтологии для управления


процессами обеспечения непрерывности и~доступности ИТ-услуг в условиях нарушения


функционирования ин\-фор\-ма\-ци\-он\-но-те\-ле\-ком\-му\-ни\-ка\-ци\-он\-ных сис\-тем


(ИТС) внутренними и~внешними негативными факторами. Приводится


базовый набор концептов, описывающих свойства системных компонентов ИТС, процессов,


ситуаций, сценариев и~отношений между ними при восстановлении штатного режима


поддержки ИТ-услуг.}





  \KW{онтология; информационно-телекоммуникационная сис\-те\-ма (ИТС); непрерывность


функционирования и~доступности сервисов (НФДС); системный подход;


предметная область; концепты}





\DOI{10.14357\!/\!08696527140414}








      \vspace*{-6pt}








\vskip 10pt plus 9pt minus 6pt








      \thispagestyle{headings}











\section{Введение}





\vspace*{-2pt}





  В последние годы все большее развитие и применение, особенно в


государственных и крупных биз\-нес-орга\-ни\-за\-ци\-ях, стали получать


современные крупномасштабные и~территориально распределенные ИТС,


интегрирующие в~себя информационные и~телекоммуникационные


составляющие различных систем и~сетей. Внедрение ИТС существенно


повысило эффективность деятельности крупных организаций


и~государственных структур. Существующая ИТС~--- это сложная


  ор\-га\-ни\-за\-ци\-он\-но-тех\-ни\-че\-ская система, представляющая собой


интегрированную пространственно распределенную автоматизированную


мультисервисную систему, предназначенную для обеспечения деятельности


организации и~предоставления должностным лицам и~структурным


подразделениям информационных и~телекоммуникационных услуг.


Непрерывность процессов управления любой автоматизированной


организацией стала существенно зависеть от устойчивости функционирования


ИТС в~различных условиях обстановки, особенно при возникновении


различного рода природных, техногенных аварий и~катастроф, различного


рода социальных катаклизмов и~интеллектуальных воздействий. Структурная


сложность и~территориальная распределенность этих систем привели


к~повышению их уязвимости от постоянно возрастающих опасных воз-\linebreak\vspace*{-12pt}





\pagebreak





\noindent


действий


окружающей среды. Увеличение сложности привело к~повышению рисков,


по\-рож\-да\-емых нарушением функционирования ИТС. Согласно одному из


последних исследований в~области бизнеса, проведенных компанией Forrester


и~результаты которых опубликованы в~докладе Forrester Report Sun Gard


2~сентября 2010~г.~[1], планирование и управление непрерывностью бизнеса


(BCP~--- Business Continuity Planning и~BCM~--- Business Continuity


Management) и~обеспечение восстановления деятельности в~непредвиденных


обстоятельствах и~бедстви\-ях (Disaster Recovery Planning~--- DRP)~--- эти два


направления деятельности определились как самые приоритетные, которые


занимают умы руководителей ведущих финансовых


  и~ИТ-ор\-га\-ни\-за\-ций на всех континентах. Эти исследования


проводились в~1-м и~2-м кварталах 2010~г.\ не только в России, но


и~в~Австралии, Новой Зеландии, Бразилии, Мексике, Канаде, США, Гонконге,


Китае, Индии, Японии, Великобритании, Германии и~Франции.





  Проблема обеспечения непрерывности автоматизированного управ\-ле\-ния


является сложным многодисциплинарным направлением работ,


обеспечи\-ва\-ющим бесперебойное функционирование крупных организаций


и~государственных учреждений, особенно при возникновении крупных


катастроф и~чрезвычайных ситуаций. Поэтому для решения этой проблемы


возникла необходимость обеспечить защиту инфраструктуры ИТС,


технологических процессов, информации, персонала, решить задачу


оперативной поддержки ИТ-услуг, восстановления нарушенных процессов


и~таким образом обеспечить устойчивую деятельность организации, защитить


ее репутацию на рынке, в~международной и~национальной сферах.





  Проектирование, создание и эксплуатация ИТС изначально ориентировались


на ее функциональность и~надежность. В~связи с~превращением


государственных и~особенно коммерческих организаций, связанных


с~бизнесом, в~автоматизированные объекты возникла проблема рассматривать


ИТС как техническое средство сервисного обслуживания, как источник


предоставления информационных и~телекоммуникационных услуг, т.\,е.\


возникла проблема непрерывно поддерживать сервисное обслуживание


деловых и~биз\-нес-про\-цес\-сов в~любых условиях, в~том числе при


чрезвычайных обстоятельствах. Смена парадигмы и~акцента применения ИТС


привели к~тому, что состояние системы и~качество предоставляемых сервисов


и~услуг стало сильно влиять на биз\-нес-про\-цес\-сы. В~то же время


требования бизнеса к~непрерывности стали сказываться на архитектуре


функционирования ИТС и~технологии их проектирования. Осознание самой


проблемы обеспечения непрерывности деятельности организаций на базе ИТС


с~применением сервисных и~процессных технологий еще только начинает


формироваться. Хотя есть стандарты, есть определенные требования


к~качеству и~устойчивости функционирования бизнеса (в~част\-ности


стандарты серии BS25999), но это в~основном организационные требования,


слабо связанные со спецификой информационных технологий (ИТ). Единых


методологий и~методик оценки устойчивости деятельности организаций,


функционирующих на базе ИТС, практически нет, хотя во многих странах


мира на законодательном уровне закреплены требования по управ\-ле\-нию


непрерывностью деятельности банковской сферы, определена структура


и~содержание плана действий и~восстановления данного вида деятельности.





  Обеспечение непрерывности функционирования и~доступности ИТ-услуг


(НФДС) в~такой сложной системе, как ИТС, в~условиях постоянно


возрастающих воздействий и~чрезвычайных ситуаций может быть достигнуто


только на основе системного подхода при проектировании ИТС, сочетающего


в~себе функциональные решения со знаниями обеспечения НФДС на основе


общих ап\-па\-рат\-но-про\-грам\-мных, информационных и~организационных


решений с~использованием единого терминологического подхода. Ключевым


моментом создания такой системы является тео\-рия и~методология единого


технологического и~организационного предметного поля, в~котором


функционируют объекты ИТС и~субъекты обеспечения НФДС. Внедрение


ИТС в~различных организациях показывает, что в~современных условиях


необходим качественно новый подход к~оценке устойчивости ИТС, переход


к~восприятию ИТС как технической базы данных и~знаний, ориентированных


на формирование и~предоставление информационных


и~телекоммуникационных услуг конечным пользователям. Ориентация на


максимальное удовлетворение потребностей конечных пользователей


обу\-слов\-ливает необходимость поднять <<планку>> технической поддержки


качества услуг путем введения в~процессы проектирования ИТС требований по


обеспечению непрерывности функционирования и~доступности


  ИТ-сер\-ви\-сов и~ИТ-услуг в~различных условиях эксплуатации. Идеология


такого подхода в~полной мере еще не сложилась. Методики единого


онтологического обеспечения высококачественных ИТ-услуг,


соответствующих требованиям пользователей ИТС в~соответствии


с~требованиями международных и~национальных стандартов управ\-ле\-ния


качеством ИТ-услуг с~учетом требований их непрерывности и~доступности


в~различных условиях обстановки, не отработаны. Возникает потребность


проведения исследований по изысканию путей более рациональных


онтологических подходов к~описанию ре\-сурс\-но-сер\-вис\-ных


возможностей ИТС для формирования и~предоставления качественных


  ИТ-услуг, соответствующих современным требованиям управ\-ле\-ния


непрерывностью бизнеса и~доступности ИТ-услуг. 





Имеющиеся пуб\-ли\-ка\-ции,


как правило, сосредоточивают свое внимание на частных вопросах и~не


рассматривают эту проблему в~комплексе научных и~технических работ,


уделяя в~основном внимание разработке общих вопросов онтологии без учета


их восприятия и~реализации знаний в~мультисервисных ИТС~[2, 3].


Необходим качественно новый взгляд на онтологию НФДС мультисервисных


ИТС. Данный подход должен базироваться на общесистемной методологии


типового применения онтологического инжиниринга и~реинжиниринга ИТС


в~де\-я\-тель\-ности организаций, на проектировании всех основных частей


и~системы в~целом с~учетом требований устойчивости, качества


  ИТ-сер\-ви\-сов, а~так\-же до\-ступ\-ности ИТ-услуг на базе принципов


ор\-га\-ни\-за\-ци\-он\-но-тех\-ни\-че\-ской интеграции знаний предметной области (ПО)


с~информационными, техническими и~программными возможностями ИТС на


базе современной технологии автоматизированных процессов управ\-ле\-ния


деятельностью организации. Для успешного решения данной проблемы


необходимо разработать тео\-рию и~методологию обеспечения непрерывности


функционирования гетерогенных автоматизированных систем органов


государственного управ\-ле\-ния и~бизнеса в~едином информационном


пространстве на базе методов компьютерной лингвистики и~онтологического


моделирования. Оставив за рамками данной статьи философские


и~тео\-ре\-ти\-че\-ские вопросы трактовки понятия онтология, рассмотрим вопросы


ее перемещения в~актуальную на сегодняшний день среду обеспечения


непрерывности функционирования государственных и~коммерческих


автоматизированных организаций на базе ИТС.





  Онтология ИТС~--- целостная структурная спецификация компонентов ПО


обеспечения бесперебойного предоставления ИТ-услуг, ее формализованное


представление, которое включает описание терминов ПО и~логических


выражений, описывающих, как они соотносятся между собой, между областями


и~внутри каждой ПО. Онтология формирует самое общее представление об


извлечении, формировании, структурировании и~предоставлении информации


(данных и знаний) о~состоянии ИТС, процессах обеспечения непрерывности


функционирования и предоставления ИТ-услуг конечному пользователю.


В~данной статье авторы попытались подойти к~решению этой проблемы


с~технической точки зрения путем разработки некоторых исходных положений


общесистемной методологии проектирования комплексной интегрированной


онтологической модели обеспечения НФДС крупномасштабной


и~территориально распределенной системы, содержащей множество


информационных и~технических ПО, с~учетом знаний поддержки семантики


ИТ-услуг. В~част\-ности, рассматриваются ресурсно-сервисный


  и~про\-цес\-сно-сис\-тем\-ный подходы с учетом распространенных в мире


стандартов управ\-ле\-ния непрерывностью ИТ-сервисов и доступностью


  ИТ-услуг~[4].


  


\vspace*{-2pt}





\section{Характеристика предметной области информационно-телекоммуникационной системы}





\vspace*{-1pt}





  Информационно-телекоммуникационная система


   является масштабной технической пространственно распределенной


системой со сложными взаимосвязями ре\-сурс\-но-сер\-вис\-ных возможностей


компонентов инфраструктуры между собой и~с~пользователями. Вследствие


этого восстановление нарушенного процесса целевого функционирования ИТС


в~случае возникновения негативных факторов становится зачастую сложной


системной проб\-ле\-мой. Исследования показывают, что ее решение известными


методами тео\-рии надежности, т.\,е.\ только устранением локальных инцидентов


в~технических и~программных средствах, не представляется возможным.


Методы тео\-рии надежности не гарантируют безотказной работы систем (сетей),


подсистем на всем технологическом цикле формирования и~предоставления


ИТ-услуг. Если в~ка\-кой-то мере методы надежности еще обеспечивают


решение проблемы отказоустойчивости, то проблему катастрофоустойчивости


они не решают.





  В этой связи весьма актуальной становится задача решения рассмат\-ри\-ва\-емой


проблемы путем создания ИТС высокой доступности с~централизованной


сис\-те\-мой управ\-ле\-ния всем процессом обеспечения НФДС. Создание


комплексной сис\-те\-мы управ\-ле\-ния НФДС обеспечивает устойчивость


организации за счет пред\-от\-вра\-ще\-ния действия негативных факторов


и,~в~случае возникновения последних, оперативного восстановления


  ИТ-сер\-ви\-сов. Система управ\-ле\-ния НФДС представляет собой целостный


процесс, в~рамках которого идентифицируются потенциальные угрозы,


оцениваются возможные риски, порождаемые ИТС, и~их воздействие на


деятельность организации, эффективно осуществляется процесс


восстановления, а~так\-же гарантируется заинтересованность обеспечивающих


сторон и~сохранение репутации организации. Системе управ\-ле\-ния НФДС


приходится иметь дело с~разнородными потоками структурированных


и~неструктурированных данных, описывающих как заранее определенные, так


и~непредсказуемые ситуации. Большая разнородность субъектов и~объектов


управ\-ле\-ния, а~так\-же непредсказуемость действия негативных факторов


требуют применения таких современных методов управ\-ле\-ния


ИТ-ор\-га\-ни\-за\-ци\-ей, которые позволяют организовать эффективно


управ\-ля\-емое и~контролируемое информационное пространство в~рамках всей


организации на базе единого взаимопонимания подразделений информатизации


и~служб поддержки ре\-сурс\-но-сер\-вис\-но\-го обеспечения пользователей


в~различных условиях обстановки. Ориентация на объединение в~едином


информационном пространстве пользователей ре\-сурс\-но-сер\-вис\-ных


процессов и~процессов управ\-ле\-ния НФДС обусловливает необходимость


решения проблемы <<семантической интероперабельности>> данных и~знаний


и~их соответствие целевой функции системы. В~европейской программе


IDABC (Interoperable Delivery of European


eGovernment Services to public Administrations, Businesses and Citizens)~[5] <<интероперабельность~--- это способность информационных


и~коммуникационных технологий (ИКТ) и~биз\-нес-про\-цес\-сов, которые они


поддерживают, к~обмену данными и~к~совместному (коллективному)


использованию информации и~знаний>>. Исходя из этого общего определения,


семантическую интероперабельность ИТС будем рассматривать как некоторый


аспект взаимодействия пользователя с~системой, который позволяет


гарантировать однозначное понимание смысла формулируемой ИТ-услу\-ги


требованиям пользователя. Наряду с семантической интероперабельностью


в~решении проблемы НФДС приобрела значимость техническая


интероперабельность, которая охватывает технические системы,


обеспечивающие бесперебойность предоставления, доступность


и~безопасность ИТ-услуг и~способность интеграции различных средств между


собой, технологические процессы формирования и~поддержки ИТ-услуг


и~регламентацию действий при возникновении опасных воздействий


и~чрезвычайных ситуаций.











  При решении этих проблем центральное место занимают исследования


в~об\-ласти онтологии, так как в~онтологическом аспекте разделение


декларативных и~процедурных знаний соответствует различным типам знаний,


включая знания об ИТ-услу\-гах и~знания об их ре\-сурс\-но-сер\-вис\-ном


формировании и~технической поддержке. В~информатике под онтологией


принято понимать формальное представление знаний в~виде множества


понятий некоторой ПО, в~данном случае~--- НФДС. Онтологическое


моделирование имеет приоритетное значение, так как отражает


многоаспектность проблемы формирования ИТ-услуг и~обеспечения НФДС


и~создает предпосылки повышения эффективности применения ИТС за счет


уменьшения семантической несогласованности процессов формирования


и~поддержки НФДС путем более глубокого их описания множеством


примитивов представления. В~част\-ности, в~контексте обеспечения НФДС


онтологию можно рассматривать как средство моделирования знаний об


индивидах про\-цес\-сно-сис\-тем\-но\-го подхода, их атрибутах и~их связях


с~другими индивидами, описывающих среду функционирования ИТС. Важным


понятием в~онтологии является ПО, которую определяют как множество


связанных различными отношениями сущностей, составляющих материальное


содержание области функционирования ИТС. В~качестве сущностей ПО


выступают объекты, процессы, средства, явления и~ситуации. В~каж\-дой ПО


определяется совокупность задач, которые решаются, используя информацию,


содержащуюся в~описаниях свойств ее сущностей. Такую совокупность,


исходя из структуры и~функциональности ИТС, будем называть ПО


компонентов ИТС. На схеме (рис.~1) представлена структура ПО, связанных


с~решением задач поддержки непрерывности и~доступности ИТ-сер\-ви\-сов.


С~появлением проблемы обеспечения НФДС возникла потребность в~более


глубоком анализе информации для принятия решений при возникновении


инцидентов и~кризисных ситуаций. Потребовалось запоминать


и~целенаправленно прогнозировать знания об отношениях между сущностями


ПО внутри системы и~отношениями системы с~сущностями внешней среды.











  Онтология ПО управления НФДС включает в~себя концепты, описы\-ва\-ющие


ре\-сурс\-но-сер\-вис\-ные возможности ИТС, ИТ-услуги, процессы, негативные


факторы и~риски, порождаемые ими, ситуации и~сюжеты, связанные


с~нарушением НФДС. Структура описания каждой категории концептов


представляет собой некоторый набор атрибутов адекватных ПО состояния


компонентов ИТС в~заданной момент времени. В~пределах структуры


онтологии для каждого компонента характерно наличие своего набора


атрибутов, отражающих семантическую специфику области и~условий его


функционирования в~ИТС. Онтологические системы ИТС на всех уровнях


строятся на основе следующих принципов:


  \begin{itemize}


\item формализация концепта ПО каждой системы, процесса, т.\,е.\ описание


элементов в~строго определенных терминах, моделях, отношениях и~т.\,д.;


  \item  использование базовых терминов и~отношений, на основе которых


конст\-ру\-иру\-ют\-ся все остальные понятия онтологий всех уровней;


  \item  полнота, внутреннее единство, логическая взаимосвязь


и~непротиворечивость.


  \end{itemize}











  Онтологические языки, создаваемые и применяемые в~ИТС для описания


ИТ-услуг, внешних и~внутренних угроз, состояния ре\-сурс\-но-сер\-вис\-ных


возможностей компонентов, относительно <<просты>> и~строятся, как


правило, на декларативных и~процедурных составляющих. К~декларативной


части относятся средства представления понятий, связей между понятиями


в~виде некоторого машинного отологического словаря. Онтологический


словарь ИТС~--- это способ формального представления знаний в~виде


множества понятий и~отношений между ними~[6].


  \begin{figure} %fig1


  \vspace*{1pt}


 \begin{center}


 \mbox{%


 \epsfxsize=124.702mm


 \epsfbox{bys-1.eps}


 }


 \end{center}


 \vspace*{-9pt}


\Caption{Классификация онтологий ИТС}


  \end{figure}


  К~процедурной части


относятся средства придания смысла единицам онтологического языка.


Онтологические языки управления НФДС условно можно разделить на два


вида: информационные и~технические языки. Информационные языки служат


для предоставления данных и~восприятия ИТ-услуг пользователями.


Технические языки служат для представления данных о~ситуациях и~сюжетах,


сложившихся в~процессе технического функционирования ИТС в~штатном


режиме и~чрезвычайных условиях на всем жизненном цикле системы. Оба этих


языка строятся на базе лексики и~грамматики естественных языков,


ориентированных на моделирование процессов понимания смысла текста


предоставляемых ИТ-услуг, формализации технических ситуаций и~сюжетов,


возникающих в~ИТС. Важной методологической проблемой при управлении


НФДС, сочетающей в~себе че\-ло\-ве\-ко-ма\-шин\-ный, системный


и~процессуальные факторы, является правильная оценка необходимого


соотношения между декларативными и~процедурными компонентами системы


анализа и~оценки состояния рисков, порождаемых нарушением


функционирования ИТС и~возможностью оперативного восстановления


штатного режима. В~вопросах соотношения между методами компьютерной


лингвистики, ориентированной на моделирование процесса понимания смысла


текста, методами онтологического моделирования, ориентированного на


представление знаний, и~методами концептуального моделирования, где на


первый план выходит описание структуры ПО с~описаниями поведения,


задаваемыми в~виде функций и~процессов, к~сожалению, ясности на


сегодняшний день нет. Процесс управления НФДС~--- это сложная


  че\-ло\-ве\-ко-ма\-шин\-ная ПО, функционирующая в~гетерогенной среде


с~внутренними и~внешними негативными воздействиями. По мнению авторов,


это обстоятельство диктует необходимость разумного сочетания системного,


процессного, ре\-сурс\-но-сер\-вис\-но\-го и~онтологического подходов, так как


при этом меньше встретится тупиков и~труднопреодолимых препятствий. Это


позволит перейти к~созданию ИТС высокой доступности с~комплексной


автоматизированной системой управления НФДС на всем жизненном цикле


ИТС в~различных условиях обстановки.





  Извлечение знаний из баз данных и~баз знаний (БЗ) основывается на


понимании семантики ПО. В~качестве первичного источника семантического


знания в~системах обычно используются машинные онтологические словари


и~классификаторы~[7]. Машинный онтологический словарь и~классификаторы


ИТС содержат множество понятий ПО функциональной и~обеспечивающей


части, поддержки ИТ-услуг, ситуации и~сюжетов. Каждая ПО характеризуется


своим специфическим набором понятий, рассматриваемых с~точки зрения ее


смысла. Как правило, система функционирования ИТС строится на базе


единого словаря и~классификатора, в~которых интегрируются


непересекающиеся понятия различных ПО. Словарь (особенно словарь типа


тезауруса) является средством устранения семантической неоднозначности


между компонентами системы, средством семантической интерпретации ПО


в~ходе общения специалистов служб поддержки НФДС между собой


и~пользователями ИТ-услуг. 





Сложность составления единого словаря


поддержки ИТ-услуг в~условиях неопределенности внутренних и~особенно


внешних факторов заключается в~том, что описать все необходимые знания


разнородных ПО в~одной системе понятий словаря и~отношений между ними


в~условиях структурированных и~неструктурированных данных,


используемых для извлечения знаний, не всегда возможно. Вторая сложность


заключается в~том, что в~каж\-дой ПО используются свои специальные языки,


средства обработки и~извлечения знаний (алгоритмические языки, языки


представления знаний, языки программной среды и~т.\,п.), а~так\-же различные


подходы и~инструментальные средства онтологической поддержки процессов


принятия решений в~условиях рисков, порождаемых неопределенностью


факторов воздействия на состояние ИТС.





\section{Структура онтологии информационно-телекоммуникационной системы}





  Для решения вопросов применения онтологии для обеспечения НФДС


рассмотрим несколько рабочих базовых определений, благодаря которым


становится возможным практическое применение онтологии в~ИТС. Это,


прежде всего, знания и~данные. В~свое время Лейбниц предложил такую


формулу: знания есть данные плюс их семантическая интерпретация. Эта


формула четко определяет связь между знаниями, исходными данными


и~средствами их смысловой интерпретации. Действительно, в~ИТС данные~---


это отдельные формализованные факторы, характеризующие объекты,


процессы, явления и~их атрибутику в~ПО, обслуживаемой системой. При


обработке в~системе они могут представляться в~структурированном или


неструктурированном виде. Хотя данные и~несут некоторую семантическую


нагрузку и~могут восприниматься как декларативные знания (скорее,


метазнания), однако каждым субъектом, работающим с~ними (в~том числе


разработчиками системы/программ), воспринимаются и~трактуются по-своему.





  Данные становятся знаниями только после того, как в~их содержание


заложены инструменты трактовки однозначности смыслового восприятия на


основе закономерностей ПО с~учетом принципов, связей и~законов,


существующих между данными и~отраженных в~их идентификаторах или


результатах представления субъектам. Иногда такие знания называют


<<хорошо структурированные данные>>. Таким образом, информационные


компоненты ИТС, включающие в~свой состав методы и~средства сбора,


хранения и~обработки информации, основанные на концепции баз и~банков


данных, где перечисленные процессы обработки ориентированы на восприятие


смысла конкретным пользователям и~не имеют развитых средств


семантической интерпретации, относятся к~классу АИС (автоматизированных


информационных систем) обработки данных. Автоматизированные информационные


сис\-те\-мы, обладающие наряду


с~процессами обработки данных инструментарием семантической


интерпретации данных и~способностью на основе более полного отображения


отношений и~связей между ними создавать и~прогнозировать новые знания


о~сущностях ПО, получили название интеллектуальных АИС.


В~таких системах структурированная модель ПО


получила название БЗ, а~инструментальные средства,


предназначенные для решения задач сохранения известных знаний


(метазнаний), синтеза новых знаний и~управления этими знаниями, приобрели


характер интеллектуальных средств. Здесь важно заметить, что концептуальной


основой понятия <<знания>> является то, что оно возникает в результате


продуктивного осмысления декларативной информации (правил, процедурных


знаний) о~сущностях ПО.





  Для современных автоматизированных государственных и~коммерческих


структур характерно постоянное изменение среды функционирования,


обуслов\-лен\-ное наличием большого числа негативных факторов, действие


которых в~первую очередь отражается на устойчивости технической


компоненты организации~--- АИС


(см.\ рис.~1). Устойчивость ИТС напрямую зависит от наличия знаний о~процессах


функционирования компонентов системы и~их спо\-соб\-ности обеспечивать


непрерывное предоставление ИТ-услуг, доступное конечному пользователю.


Онтологический подход к~проектированию системы управления знаниями


позволяет подразделениям информатизации осуществлять целостный,


системный подход к~ре\-сурс\-но-сер\-вис\-но\-му анализу ПО компонент ИТС


и~оценки их состояния, что позволяет восстанавливать нарушенные логические


связи между атрибутами ПО.





  На формальном уровне онтология представляет собой набор концептов


(понятий), свойств концептов (атрибутов, ролей), отношений между


концептами (зависимостей, функций) и~дополнительных ограничений.


Концептами в ИТС являются информационные задачи, функции систем (сетей)


и~подсистем, технологические процессы сбора, процессы обработки данных


о~состоянии процессов формирования ИТ-услуг, действия по устранению


возникших нарушений в~процессах НФДС и~т.\,п.





  Классифицировать онтологию в ИТС можно по различным признакам


(основаниям) в~зависимости от стоящей цели. Например, по прикладным ПО:


область описания содержания ИТ-услуг, область содержания процессов


формирования информационных и~телекоммуникационных ИТ-услуг, область


управления НФДС, область оценки угроз и~рисков и~т.\,д. По уровню


прикладных ПО, исходя из инфраструктуры ИТС и~процессов решения задач


сервисного обслуживания конечных пользователей, целесообразно различать


следующие уровни онтологии (см.\ рис.~1).





  \textit{Верхний уровень: <<онтология предметной области организации>>.}


Описывает наиболее общие концепты, которые ориентированы как на ПО


организации, так и~на концепты прикладных областей инфраструктуры ИТС.


Эта онтология предназначена для совместного использования концептов


и~отношений между ними как в организации, так и~в~ИТС. Основой этой


онтологии является БЗ, содержащая все понятия организации и~ИТС


в~виде единого отраслевого стандартного словаря. Например, системная


онтология ИТС Банка России строится на базе финансовых терминов


и~базовых терминов, описывающих область функционирования ИТС.





  \textit{Онтологии ИТ-услуг.} Эти онтологии используются как при описании


выходной функциональности ИТ-услуг, так и~при описании требований


к~обеспечивающим их приложениям, предназначенным для формирования


конкретных ИТ-услуг, поддержки пользователей, а~так\-же обеспечения


информационной и~операционной безопасности процессов деятельности


организации. Данные онтологии отражают специфику


  ре\-сурс\-но-сер\-вис\-ных процессов формирования ИТ-услуг, отношения


между концептами угроз и~их взаимосвязь с~концептами уязвимостей ИТС,


безопасностью, порождаемыми ИТС рисками и~т.\,д.





  \textit{Онтология уровня приложений.} Онтологии этого уровня


ориентированы на использование при описании конкретных задач как в составе


процессов формирования ИТ-услуг, так и~в~процессах поддержки


пользователей, а~так\-же в~процессах обеспечения информационной


и~операционной безопасности. Одновременно онтологии данного уровня


содержат специальные термины, присущие описанию процессов в~аппаратных


и~программных средствах систем.





  \textit{Уровень системных онтологий.} Онтологии этого уровня включают в~себя


онтологии функциональных, обеспечивающих систем, онтологию средств


обеспечения устойчивости и~безопасности ИТС. Примером такой онтологии


может служить БЗ, содержащая понятия функциональных


информационных систем, телекоммуникационных систем, систем


  опе\-ра\-тив\-но-тех\-ни\-че\-ско\-го управ\-ле\-ния ИТС и~т.\,п.





  Одним из основных требований к онтологиям всех уровней является наличие


терминологической и~понятийной взаимосвязи между ними. Строение


и~свойства всех процессов непрерывного функционирования любой


ИТС, равно как и~процесс решения той или иной прикладной задачи, могут


быть эффективно исследованы при помощи словаря терминов, адекватно


описывающего характеристики объектов, систем и~процессов решения, при


наличии точных и~однозначных определений терминов и их классификации


в~виде логической иерархии.





  Каждому онтологическому уровню свойствен определенный набор


атрибутов. К~ним относятся:


  \begin{enumerate}[1.]


  \item \textbf{Концепты-системы}~--- объекты, которые описываются


списком собственных атрибутов, характеризующих структуру ПО,


функциональность и~критическую важность ре\-сурс\-но-сер\-вис\-ных


возможностей ИТС.


  \item \textbf{Концепты-объекты} описываются именами компонент ИТС,


каждое из которых имеет свои атрибуты, описывающие состав, структуру


и~функциональность. Описание аппаратной или программной компоненты


ИТС соответствует кон\-цеп\-ту-свой\-ст\-ву с~отражением своей


специфической ПО. Кон\-цеп\-там-объ\-ек\-там в~терминологии Resource


Description Framework (RDF) соответствуют концепты предметов~--- субъектов


и~объектов.


  \item \textbf{Концепт-свойство} характеризует состояние устойчивости


систем и~объектов к~негативным воздействиям (например, свойства отказо-


и~катастрофоустойчивости). Кон\-цеп\-та\-ми-свой\-ст\-ва\-ми могут быть типы


данных, описывающих значения и~методы оценки устойчивости и~рисков,


которые позволяют судить о~мере близости понятий при выработке решений


на основе логического вывода на неточных и~неполных данных. В~RDF


  кон\-цеп\-там-свой\-ст\-вам соответствуют предикаты.


  \item \textbf{Концепты-отношения} используются для отображения между


понятиями ПО парадигматических и~синтагматических отношений, которые


отражают связь как между понятиями внутри ПО, так и~между понятиями


различных ПО.


  \item Концепты-процессы характеризуют преобразование


  ре\-сурс\-но-сер\-вис\-ных возможностей ИТС в ИТ-услуги. С~точки зрения


управления НФДС в~категорию <<процессы>> попадают понятия, атрибуты


которых описывают требования бенефицианта\!/\!реципиента ИТ-услуги,


состояние негативных факторов, сценарий их воздействия на компоненты ИТС,


каналы воздействия, результаты нарушения функциональности и~доступности


сервисов и~услуг, риски и~т.\,д.


  \item \textbf{Ситуация}~--- совокупность значений атрибутов,


характеризующих состояние функционирования ИТС, доступность ИТ-услуг,


непрерывность процессов их формирования, реакцию на внутренние и~внешние


воздействия и~отношения между атрибутами в~некоторый момент времени.


  \item \textbf{Сюжет}~--- упорядоченная во времени последовательность


ситуаций, име\-ющих системную и~процессную пересекающуюся взаимосвязь.


Сюжет может быть представлен в~виде сценария, сетевого графика, диаграммы


Ганта для описания плана обеспечения НФДС и~проведения работ по


восстановлению штатного режима деятельности (плана ОНиВД).


  \end{enumerate}





  Управление НФДС связано с~перечисленными концептами в~процессе


построения запросов, получения пользователем ИТ-услуг, а~так\-же


в~процессе эксплуатации ИТС как с~формой представления знаний


о~соответствующей ПО. В~качестве первичного источника семантических


знаний для каждой ПО и~описания процессов используется свой


специфический онтологический язык. Таким образом, ИТ-услу\-ги, процессы


их формирования, предоставления и~поддержки описываются своим


специфическим набором онтологий, свойственных функциональному


предназначению и~техническому построению ИТС, которые позволяют


рассматривать процессы обеспечения НФДС с~точки зрения их смысла.


Устойчивое функционирование ИТС может быть достигнуто только созданием


эффективной системы управления знаниями на основе автоматизированной


семантической обработки информации для принятия решений по обеспечению


НФДС и~для выполнения принятого решения. Управление знаниями


о~НФДС~--- совокупность процессов и~технологий, предназначенных для


выявления негативных факторов, создания, обработки и~хранения данных


и~ресурсов для устранения последствий их воздействия на ИТ-услу\-ги,


предоставляемые конечному пользователю. Для решения задач


централизованного управления НФДС потребуется более сложная


  ло\-ги\-ко-се\-ман\-ти\-че\-ская обработка информации о~сущностях каждой


ПО, улучшения структуры информационного обеспечения, создания


универсальных баз данных и~БЗ. Под БЗ будем понимать


составную часть каждой функциональной, обеспечивающей части ИТС,


представляющей собой комплекс языковых алгоритмических и~программных


средств, предназначенных для восприятия, обработки, хранения и~выдачи


знаний о~ПО. Этот комплекс включает:


  \begin{itemize}


  \item формализованные факты и данные, отражающие модель состояния ПО;


  \item правила и программы, позволяющие рассчитывать показатели,


эффективно и~целенаправленно логически преобразовывать данные в~ходе


анализа и~синтеза при решении задач оценки состояния ПО и~выработки


управляющих воздействий по устранению языковой неоднозначности между


компонентами ИТС.


  \end{itemize}





  Онтологическая модель позволяет формализовать знания о~ПО каждой


сис\-те\-мы, ее состояний и~процессах управления, которые являются


центральным концептом, с~которым связано применение основных понятий,


определенных в~международных и~национальных стандартах ISO


(International Organization for Standartization)~[8--11].


Применение онтологии, описывающей как процессы обеспечения НФДС, так


и~процессы функционирования ИТС в~целом, поз\-во\-ля\-ет производить


построение достоверных утверждений о~состоянии отказо-


и~катастрофоустойчивости сис\-те\-мы в~отношении предоставления ИТ-услуг


в~требуемые моменты времени. На основе этих утверждений могут быть


приняты решения, позволяющие поддерживать критически важные ИТ-услуги


в~кризисных ситуациях, а~так\-же проводить работы по восстановлению


штатного режима функционирования ИТС. Важно учитывать то


обстоятельство, что онтология является не только описанием ПО, но


и~методологией семантического описания характеристик, связанных


с~исследованиями функционального и~технического состояния устойчивости


на основе знаний о~кризисных ситуациях, требующих оперативного принятия


решений по обеспечению НФДС.





  Выбор онтологической модели описания процессов обеспечения НФДС


обуслов\-лен тем, что этот подход позволяет на ло\-ги\-ко-се\-ман\-ти\-че\-ском


уровне объединить в~едином представлении ИТ-услу\-ги,


  ре\-сурс\-но-сер\-вис\-ную модель, системные процессы, модели угроз


и~построить единую модель обеспечения НФДС на базе единого


онтологического языка. Онтологическая модель ИТС позволяет создавать


и~вести БЗ системы управления НФДС в~кризисных ситуациях и~выполнять


следующие функции:


  \begin{itemize}


  \item  оценивать ИТ-услуги на семантическое соответствие требованиям


конечного пользователя;


  \item распознавать состояние технологических процессов


  ре\-сурс\-но-сер\-вис\-но\-го формирования ИТ-услуг;


  \item  распознавать кризисные ситуации, нештатные ситуации (НШС)


и~инциденты, нарушающие НФДС;


  \item  прогнозировать влияние кризисных ситуаций, НШС, инцидентов на


НФДС;


  \item  планировать разработку единых планов ликвидации последствий


кризисных ситуаций и~устранять локальные и~системные НШС и~инциденты;


  \item  оценивать риски, порождаемые ИТС, и~планировать мероприятия по


их снижению на объектах и~системах ИТС;


  \item  планировать мероприятия по повышению качества и~готовности


превентивных мер к~противодействию угрозам.


  \end{itemize}





  Многие проблемы, порождаемые негативными факторами и~их


последствиями (инцидентами, авариями, катастрофами и~другими НШС),


приводят к~нарушению НФДС и~требуют наличия описания в~виде моделей


и~программных комплексов, реализующих последние. Для интеграции


различных моделей обеспечения НФДС с~описаниями процессов


функционирования ИТС предлагается использовать БЗ, опирающуюся


на онтологическую модель НФДС. Это означает, что состояние каждой


функциональной обеспечивающей системы, которая позиционируется


моделями обеспечения устойчивости, соотносится с~БЗ, в~рамках


которой рассматривается ло\-ги\-ко-се\-ман\-ти\-че\-ская связь между


компонентами ИТС, угрозами и~моделями ситуаций. Концепт устойчивости


детализируется в~схему системы управления НФДС, которая может включать


концепт таких моделей, как ло\-ги\-ко-се\-ман\-ти\-че\-ское формирование


сервисов и~ИТ-услуг, описание угроз, оценка критичности объектов, систем


и~сервисов ИТС, обеспечение отказоустойчивости технических


и~программных средств, поддержка катастрофоустойчивости, оценка рисков,


восстановление функционирования и~доступности ИТ-услуг, обеспечение


безопасности, организация эксплуатации ИТС в~повседневных


и~чрезвычайных условиях и~т.\,п. Такая схема обычно представляет собою


иерархическую модель данных. Эти данные должны обеспечить возможность


реализации работ по выполнению планов ОНиВД организаций.


{\looseness=1





}





\section{Онтологическая модель обеспечения безопасности


информационно-телекоммуникационной системы}





  Как уже было сказано выше, обеспечение НФДС не является самоцелью.


Нарушения НФДС актуальны постольку, поскольку порождают риски для


процессов основной деятельности организации. Для того чтобы установить


адекватные парадигматические связи между ПО НФДС и~ПО основной


деятельности организации, необходимо ввести в~рассмотрение особый раздел


ИТ-безопас\-ности, отвечающий за контроль и~управление связями между


нарушениями НФДС, с~одной стороны, и~порождаемыми ими~--- с~другой.





  Действительно, ИТ-безопасность в~современном ее понимании уже не может


ограничиваться проблемами информационной безопасности~[12]. Достаточно


привести примеры остановки операционных процессов крупнейших


зарубежных бирж, произошедших из-за нарушений НФДС. <<В~апреле


2000~г.\ на 8~ч (почти на всю торговую сессию) были парализованы


торги на Лондонской фондовой бирже (LSE) из-за ошибок в~работе


программного обеспечения. Убытки оценили в~несколько миллионов фунтов.


В~сентябре 2008~г.\ там же был зафиксирован сбой и~работа биржи была


остановлена на 4~ч. Это не только привело к~прямым потерям~---


около 40\% дохода биржи приходится на продажу информации о~курсах акций


в~режиме реального времени, но и~пошло на пользу конкурентам LSE.


Известны также случаи крупных аварий, ставших причиной остановки


критически важных сервисов на Токийской фондовой бирже и~NASDAQ.


Очевидно, что обеспечение непрерывности биз\-нес-про\-цес\-сов (business


continuity) является первостепенной задачей финансовых учреждений>>~[13].


Становится все более очевидным, что бесперебойность функционирования


в~важнейших отраслях современной деятельности (финансовой, техногенной,


оборонной и~др.)\ обусловливается непрерывностью функционирования


и~доступностью ИТ-сер\-ви\-сов, обслуживающих эти процессы. Роль и~место


операционной безопас\-ности в~процессе функционального автоматизирования


организации, в~отличие от информационной безопасности, авторы подробно


рассматривают в~своих работах~[14, 15]. Понятие <<операционная


безопасность>> определяется в~соответствии с~ГОСТ Р~1.0-92~[16] как


<<отсутствие недопустимого риска нанесения пользователям сервисов ИТС


ущерба вследствие нарушения НФДС>>. Онтология операционной


безопасности ИТС строится исходя из содержания ПО общей безопасности


ИТС (рис.~2).


{\looseness=1





}





\begin{figure} %fig2


\vspace*{1pt}


 \begin{center}


 \mbox{%


 \epsfxsize=125.335mm


 \epsfbox{bys-2.eps}


 }


 \end{center}


 \vspace*{-9pt}


\Caption{Содержание элементов структуры общей безопасности функционирования ИТС}


\end{figure}





  Система операционной безопасности ИТС, в~свою очередь, призвана


обеспечивать устойчивость основных процессов деятельности организации


путем выявления наиболее рисковых внутренних процессов ИТС,


предотвращения действия негативных факторов и~оперативного парирования


в~случае их возникновения. Решение задач реактивного и~проактивного


обеспечения операционной безопасности управления рисками требует


использования моделей угроз, ре\-сурс\-но-функ\-ци\-о\-наль\-ных моделей,


конфигурационных баз данных, явля\-ющих\-ся <<зачатками>> корпоративной


онтологической модели. Указанные модели являются ключевой основой


<<Центра управления системой>> (Service Desk) (рис.~3).


{\looseness=1





}








  Таким образом, в рамках системы обеспечения НФДС операционная


безопасность становится важной компонентой целостного процесса, в~рамках


которого сначала (в~рамках системы управления операционным риском)


идентифицируются потенциальные угрозы, оцениваются возможные риски,


порождаемые нарушением НФДС, и~их воздействие на процессы основной


деятельности организации.


\begin{figure} %fig3


\vspace*{1pt}


 \begin{center}


 \mbox{%


 \epsfxsize=125.396mm


 \epsfbox{bys-3.eps}


 }


 \end{center}


 \vspace*{-11pt}


\Caption{Место операционной безопасности в структуре ИТ-деятельности}


\end{figure}


Затем (в рамках уже системы обеспечения НФДС),


в~соответствии с~разработанными стратегиями обеспечения НФДС,


задействуются необходимые механизмы противодействия критичным угрозам


с~целью возобновления НФДС и~(в~дальнейшем) восстановления нарушенных


ресурсов ИТС (рис.~4).





\vspace*{-6pt}





\section{Заключение}





\vspace*{-2pt}





  К сожалению, в настоящее время нет четких границ между методами


компьютерной лингвистики, онтологией и~системами классификации. Это


обусловлено тем, что нет четкого понимания того, как соотносятся между


собой такие понятия, как <<знания>>, <<данные>>, равно как и~средства их


семантического представления в~АИС.


\begin{figure} %fig4


\vspace*{1pt}


 \begin{center}


 \mbox{%


 \epsfxsize=128.608mm


 \epsfbox{bys-4.eps}


 }


 \end{center}


 \vspace*{-11pt}


\Caption{Парадигматические связи между понятиями, обозначающими


основные виды нештатных ситуаций}


\end{figure}


Основная особенность задачи


построения онтологической модели обеспечения НФДС заключается в~том, что


в~описании ИТС присутствуют объекты (системы, технологические процессы,


сетевые структуры, разнообразные уязвимости технических и~программных


средств, базы данных и~знаний, средства защиты от угроз и~т.\,д.) различной


семантической природы, для формального описания которых не существует


однозначно совместимой терминологии, либо последней явно недостаточно


для описания существующих объектов. Кроме того, ИТС проектируются


различными разработчиками, использующими различные термины и~их


семантическую трактовку.





  Требуется разработка и применение особой технологии построения


онтологических моделей ИТС. В~качестве основы такой технологии может


быть использован про\-цес\-сно-сис\-тем\-ный подход~[4]. Основная идея


<<процессного системного подхода>> состоит в~том, что любая система


рассматривается как совокупность устойчиво повторяемых классов процессов


и~обеспечивающих их ресурсов, в~том числе элементов структуры самой


системы и~внешних ресурсов (<<входов процессов>>). Применение


<<процессного системного подхода>> к~разработке тематической онтологии


предусматривает, что


  \begin{itemize}


  \item все понятия в онтологии разделяются на аксиоматические


<<категории>> (число которых минимизируется) и~<<системы>> (объекты,


связи и~правила);


  \item  поскольку связи между объектами и~правила, действующие


в~онтологии, рассматриваются как классы процессов, то все они определяются


как <<действия>> или производные от <<действий>>;


  \item структуризация понятий <<по темам>> рассматривается в~привязке


к~структуре соответствующих процессов.


  \end{itemize}





  На основании сформулированных принципов представляется возможным


сделать попытку решить следующие частные задачи:


  \begin{itemize}


  \item построить алгоритм автоматизированного выявления основных понятий


соответствующей тематической области;


  \item построить алгоритм автоматизированного выявления <<действий>>


и~производных от <<действий>>;


  \item  построить алгоритм автоматизированной структуризации тематических


областей;


  \item  построить алгоритм выявления связей и~правил соответствующей


тематической области.


  \end{itemize}





{\small\frenchspacing


{%\baselineskip=10.8pt


\addcontentsline{toc}{section}{Литература}


\begin{thebibliography}{99}


\bibitem{1-bs} %1


\Au{Balaouras S.} Business continuity and disaster recovery are top IT priorities for


2010 and 2011~/\!\!/ For Security \& Risk Professionals, September 2, 2010.


  {\sf http:// www.rbzaneadvisors.com/pdf/Forrester\_bus\_cont\_disaster\_recovery.pdf}.





  \bibitem{3-bs} %2


\Au{Шемакин Ю.\,И.} Тезаурус в~автоматизированных системах управления


и~обработки информации.~--- М: Воениздат, 1974. 192~с.





\bibitem{2-bs} %3


\Au{Калиниченко Л.\,А.} Экспансия онтологий: <<онтологически


базированные>> информационные системы~/\!\!/ Онтологическое


моделирование: состояние, направления исследований и~применения: Тр.


II~симпозиума.~--- М.: ИПИ РАН, 2011. С.~29--44.





\bibitem{4-bs}


\Au{Быстров И.\,И., Радоманов С.\,И.} Конец мифа о~противоречивости


матричного управления~/\!\!/ Information Management: научно-методический


журнал, 2013. №\,7. С.~38--49.


\bibitem{5-bs}


Европейская комиссия. Программа IDABC (Interoperable Delivery of European


eGovernment Services to public Administrations, Businesses and Citizens). {\sf


http://web. archive.org/web/20060617105314/ec.europa.eu/idabc/en/chapter/3}.


\bibitem{6-bs}


\Au{Митрофанова О.\,А., Константинова Н.\,С.} Онтологии как системы


хранения знаний~/\!\!/ Ин\-фор\-ма\-ци\-он\-но-ком\-му\-ни\-ка\-ци\-он\-ные


технологии в образовании, 2008. 54~c. {\sf


http://www.ict.edu.ru/ft/005706/68352e2-st08.pdf}.


\bibitem{7-bs}


\Au{Белоногов Г.\,Г., Калинин Ю.\,П., Хорошилов~А.\,А.} Компьютерная


лингвистика и~перспективные информационные технологии.~--- М.: Русский


мир, 2004. 246~с.











\bibitem{9-bs} %8


Управление непрерывностью бизнеса. Ч.~1: Практическое руководство~/


Не\-офиц. пер. британского стандарта BS 25999-1:2006.


\bibitem{11-bs} %9


PAS 77:2006. IT Service Continuity Management. Code of practice. Publicly


available specification.~--- BSI, August~11, 2006.


\bibitem{10-bs} %10


Управление непрерывностью бизнеса. Ч.~2: Спецификация~/ Пер.


британского стандарта BS 25999-2:2007, выполненный ООО <<ГлобалТраст


Солюшинс>> с~разрешения Британского института стандартов (BSI) по


лицензии №\,2006AT0005.~--- BSI, 2007.





\bibitem{8-bs} %11


Управление непрерывностью информационных и коммуникационных


технологий~--- практические правила~/ Пер. британского стандарта BS


25777:2008, выполненный ООО <<ГлобалТраст Солюшинс>> с~разрешения


Британского института стандартов (BSI) по лицензии №\,2008JK0022.~--- BSI,


2008.


\bibitem{12-bs} %12


ГОСТ Р 53114-2008. Защита информации. Обеспечение информационной


безопасности в~организации. Основные термины и~определения.~--- М.:


Стандартинформ, 2009. 40~с.


\bibitem{13-bs}


\Au{Куперман М., Аверьянов Д.} Резервный центр обработки данных. Оценка


надежности~/\!\!/ Электроника (наука, технология, бизнес), 2010. Вып.~4. {\sf


http://www. electronics.ru/journal/article/75}.


\bibitem{14-bs}


\Au{Быстров И.\,И., Радоманов С.\,И.} Операционная безопасность


 ИТ-сис\-тем~/\!\!/ Сис\-тем\-ный анализ, управление и навигация: Сб. тезисов


докл. 17-й Международной научной конф.~--- М.: МАИ-Принт, 2012. С.~86.


\bibitem{15-bs}


\Au{Быстров И.\,И., Радоманов С.\,И., Сычёв~В.\,Н.} Операционная


безопасность ИТ-сис\-тем~/\!\!/ Стандартизация, сертификация, обеспечение


эффективности, качества и~безопасности информационных технологий


(ИТ-стандарт 2012): Мат-лы III~Междунар. конф. {\sf


http://www.myshared.ru/slide/482219}.


\bibitem{16-bs}


ГОСТ Р 1.0-92. Государственная система стандартизации Российской


Федерации. Основные положения.~--- М.: Госстандарт России, 1994. 23~с.


\end{thebibliography}


} }








\vspace*{-6pt}





\hfill{\small\textit{Поступила в редакцию 16.09.14}}





%\newpage








\vspace*{9pt}





\hrule





\vspace*{2pt}





\hrule





%\vspace*{9pt}








\def\tit{ONTOLOGY SUPPORT OF~CONTINUITY OF~IT-SERVICES IN~INFORMATION


AND~TELECOMMUNICATIONS SYSTEMS}





\def\titkol{Ontology support of continuity of IT-services in~ITS}


%information and~telecommunications systems}











\def\aut{I.\,I. Bystrov, S.\,I. Radomanov, and~V.\,N. Veselov}


\def\autkol{I.\,I. Bystrov, S.\,I. Radomanov, and~V.\,N. Veselov}








\titel{\tit}{\aut}{\autkol}{\titkol}





\vspace*{-12pt}





\noindent


Institute of Informatics Problems, Russian Academy of Sciences,


44-2 Vavilov Str., Moscow 119333, Russian


Federation











\def\leftfootline{\small{\textbf{\thepage}


\hfill Sistemy i Sredstva Informatiki~---


Systems and Means of Informatics 2014 vol~24 no\,4}


}%


 \def\rightfootline{\small{Sistemy i Sredstva Informatiki~---


 Systems and Means of Informatics   2014   vol~24   no\,4


\hfill \textbf{\thepage}}}











\Abste{The article discusses the possibility of using ontologies


to manage processes to ensure continuity of operations and accessibility


of IT-services\linebreak\vspace*{-12pt} }





\Abstend{under conditions of disruption of information and telecommunications system


(ITS) functioning by


internal and external negative factors. The article provides a basic


set of concepts that describe properties of ITS components,


processes, situations, scenarios, and relations between them when restoring


the native mode of support of it services.}





  \KWE{ontology; information and telecommunications system (ITS);


  continuity of operations and accessibility of IT-services;


  systematic approach; subject area; concepts}





\DOI{10.14357\!/\!08696527140414}





%\vspace*{-24pt}





%\Ack


%\noindent








\vspace*{-2pt}











\renewcommand{\bibname}{\protect\rmfamily References}


%\renewcommand{\bibname}{\large\protect\rm References}





{\small\frenchspacing


{%\baselineskip=10.8pt


\addcontentsline{toc}{section}{References}


\begin{thebibliography}{99}


\bibitem{1-bs-1}


\Aue{Balaouras, S.} 2010. Business continuity and disaster recovery are top IT


priorities for 2010 and 2011. Available at: {\sf


http://www.rbzaneadvisors.com/ pdf/Forrester\_ bus\_cont\_disaster\_recovery.pdf }(accessed August 29, 2014).





\bibitem{3-bs-1} %2


\Aue{Shemakin, Yu.\,I.} 1974. \textit{Tezaurus v avtomatizirovannykh


sistemakh upravleniya i~obrabotki informatsii} [Thesaurus in automated control


systems and information processing]. Moscow: Voenizdat. 192~p.





\bibitem{2-bs-1} %3


\Aue{Kalinichenko, L.\,A.} 2011. Ekspansiya ontologiy: ``Ontologicheski


bazirovannye'' informatsionnye sistemy [Expansion of ontology: ``An


ontologically based'' information systems]. \textit{Tr. II~Simpoziuma


``Ontologicheskoe Modelirovanie: Sostoyanie, Napravleniya Issledovaniy


i~Primeneniya''} [2nd Symposium on


Ontological Modeling: Status, Trends, Research and Application


Proceedings]. Moscow: IPI RAN. 29--44.








\bibitem{4-bs-1} %4


\Aue{Bystrov, I.\,I., and S.\,I.~Radomanov}. 2013. Konets mifa


o~protivorechivosti matrichnogo upravleniya [The end of inconsistency matrix


management myth]. \textit{Information Management: Nauchno-Metodicheskiy


Zhurnal} [The Information Management: Scientific-Methodical J.]


7:38--49.


\bibitem{5-bs-1}


European Commission. The IDABC programme (Interoperable


Delivery of European eGovernment Services to public Administrations,


Businesses and Citizens). Available at: {\sf


http://web.archive.org/web/20060617105314/ec.europa.eu/idabc/en/chapter/3}


(accessed August 29, 2014).


\bibitem{6-bs-1}


\Aue{Mitrofanova, O.\,A., and N.\,S.~Konstantinova}. 2008. Ontologii kak


sistemy khraneniya znaniy [Ontology as a storage system of knowledge].


\textit{Informatsionno-kommunikatsionnye tekhnologii v obrazovanii. Elektronnyy


resurs}. [Information and communication technologies in education. Electronic


resource]. 54~p. 


Available at: {\sf http://www.ict.edu.ru/ft/005706/68352e2-st08.pdf}


(accessed November~1, 2014).


\bibitem{7-bs-1}


\Aue{Belonogov, G.\,G., Yu.\,P.~Kalinin, and A.\,A.~Khoroshilov}. 2004.


\textit{Komp'yuternaya lingvistika i~perspektivnye informatsionnye tekhnologii}.


[Computational linguistics and information technologies]. Moscow: Russkiy Mir.


246~p.











\bibitem{9-bs-1} %8


BS 25999-1:2006. November 2006.


Business continuity management.


Code of practice. Replayced by BS ISO 22313:2012.





\bibitem{11-bs-1} %9


PAS 77:2006. August 2006.


IT service continuity management. Code of practice. Withdrawn.





\bibitem{10-bs-1} %10


BS 25999-2:2007.


Business continuity management. Specification. Replaced by BS EN ISO 22301:2014.





\bibitem{8-bs-1} %11


BS 25777:2008. December 2008.


Information and communications technology continuity management.


Code of practice. Replaced by BS ISO\!/\!IEC 27031:2011.





\bibitem{12-bs-1}


GOST R 53114-2008. 2009. \textit{Zashchita informatsii. Obespechenie


informatsionnoy bezopasnosti v organizatsii. Osnovnye terminy i opredeleniya}


[Protection of information. Information security in the organization. Basic terms


and definitions]. Moscow: Standartinform. 40~p.


\bibitem{13-bs-1}


\Aue{Kuperman, M., and D.~Aver'yanov}. 2010. Rezervnyy tsentr obrabotki


dannykh. Otsenka nadezhnosti [Backup data center. Assessment of reliability].


\textit{Elektronika (Nauka, Tekhnologiya, Biznes)} [J.~Electronics


(Science, Technology, Business)] 60(4). Available at: {\sf


http://www.electronics.ru/journal/article/75} (accessed August~29, 2014).


\bibitem{14-bs-1}


\Aue{Bystrov, I.\,I.,  and S.\,I. Radomanov}. 2012. Operatsionnaya bezopasnost'


IT-sistem [Operational security of IT systems]. \textit{Sb. tezisov dokl. 17-y


Mezhdunar. nauch. konf. ``Sistemnyy Analiz, Upravlenie i~Navigatsiya''}


[A~collection of abstracts of the 17th Scientific Conference (International)


``System Analysis, Control, and Navigation'']. Moscow: MAI. P.~86.


\bibitem{15-bs-1}


\Aue{Bystrov, I.\,I., S.\,I. Radomanov, and V.\,N.~Sychev}. 2012.


Operatsionnaya bezopasnost' IT-sistem [Operational security of IT- systems].


\textit{Doklad na III Mezhdunar. konf. ``Standartizatsiya, Sertifikatsiya,


Obespechenie Effektivnosti, Kachestva i~Bezopasnosti Informatsionnykh


Tekhnologiy, IT-Standart 2012''}  [Presentation at the 3rd Conference (International)


``Standardization, Certification, Ensuring Efficiency, Quality, and Security of


Information Technology, IT  Standard 2012]. Moscow. Available at: {\sf


http://www.myshared.ru/slide/482219/} (accessed August~29, 2014).


\bibitem{16-bs-1}


GOST R 1.0-2004. 2005. Standartizatsiya v Rossiyskoy Federatsii. Osnovnye


polozheniya [GOST R 1.0-2004 Standardization in the Russian Federation. Basic


provisions]. Moscow: IPK Izdatel'stvo standartov. 23~p.





\end{thebibliography}


} }








\vspace*{-6pt}





\hfill{\small\textit{Received September 16, 2014}}





\Contr





\noindent


\textbf{Bystrov Igor I.} (b.\ 1931)~--- Doctor of Science in technology, professor, Head of


Department, Institute of Informatics Problems,


Russian Academy of Sciences, 44-2 Vavilov Str., Moscow 119333, Russian Federation; ibystrov@ipiran.ru





\vspace*{3pt}





\noindent


\textbf{Veselov Vitaly N.} (b.\ 1940)~--- Candidate of Science (PhD) in technology, leading


scientist, Institute of Informatics Problems,


Russian Academy of Sciences, 44-2 Vavilov Str., Moscow 119333, Russian Federation; vveselov@ipiran.ru





\vspace*{3pt}





\noindent


\textbf{Radomanov Sergey I.} (b.\ 1953)~---


scientist, Institute of Informatics Problems,


Russian Academy of Sciences, 44-2 Vavilov Str., Moscow 119333, Russian Federation; radomanov@ipiran.ru











 \label{end\stat}





\renewcommand{\bibname}{\protect\rm Литература}





